\section{DEP-004: Lockfile Not Published}
\label{sec:dep-004}

\subsection*{Metadata}
\begin{tabular}{ll}
\textbf{Issue ID} & DEP-004 \\
\textbf{Severity} & Medium \\
\textbf{Category} & Dependencies \\
\textbf{Branch} & \branchref{fix/DEP-004-lockfile-publish} \\
\textbf{Changes} & \branchdiff{fix/DEP-004-lockfile-publish} \\
\end{tabular}

\subsection*{Executive Summary}
The \texttt{package.json} \texttt{files} field only includes \texttt{["dist", "README.md"]}, which means the \texttt{package-lock.json} is not included in the published npm package. Consumers who install \texttt{gitclaw} via \texttt{npm install} get non-deterministic dependency resolutions---npm resolves each dependency at install time based on version ranges, which may differ from the versions tested against.

Transitive dependencies are especially vulnerable: a minor/patch update in a transitive dependency could introduce a breaking change or security vulnerability without the project's knowledge. With 11 direct dependencies and hundreds of transitive dependencies using caret ranges, the risk of unreproducible installs is significant.

\subsection*{Root Cause Analysis}

The package configuration excluded the lockfile:

\begin{lstlisting}[language=json]
// package.json:21-24 — Before fix
"files": [
    "dist",
    "README.md"
],
\end{lstlisting}

The \texttt{files} field controls which files are included when the package is published to npm. The lockfile (\texttt{package-lock.json}) is excluded by this configuration. Even if it exists in the repository, it is not included in the published tarball.

All dependencies use \texttt{^} (caret) ranges, meaning npm installs any compatible minor/patch version:

\begin{itemize}
    \item \texttt{ws: "\^{}8.19.0"} could resolve to \texttt{8.19.0} through \texttt{8.20.x}
    \item \texttt{@opentelemetry/api: "\^{}1.9.0"} could resolve to \texttt{1.9.0} through \texttt{1.10.x}
    \item \texttt{baileys: "\^{}7.0.0-rc.9"} could resolve to any 7.x release
\end{itemize}

\subsection*{Fix Implementation}

The fix adds \texttt{package-lock.json} to the \texttt{files} field:

\begin{lstlisting}[language=diff]
// package.json — After fix
"files": [
    "dist",
-    "README.md"
+    "README.md",
+    "package-lock.json"
],
\end{lstlisting}

This ensures that \texttt{npm install} on the published package resolves dependencies to exactly the same versions tested against.

\subsection*{Affected Files}
\begin{itemize}
    \item \srcfile{package.json} --- added \texttt{package-lock.json} to \texttt{files} array
    \item \srcfile{src/\_\_tests\_\_/dep-004.test.ts} --- new lockfile verification tests
\end{itemize}

\subsection*{Verification}
\begin{lstlisting}[language=bash]
# dep-004-verify.sh
echo "Test 1: Lockfile exists in repository..."
if [ -f "package-lock.json" ]; then
    echo "  PASS"
else
    echo "  FAIL: package-lock.json not found"
    exit 1
fi

echo "Test 2: Lockfile listed in package.json files..."
if grep -q "package-lock.json" package.json; then
    echo "  PASS"
else
    echo "  FAIL: package-lock.json not in files field"
    exit 1
fi
\end{lstlisting}

\begin{lstlisting}[language=TypeScript]
// dep-004.test.ts
test("package-lock.json exists in project root", () => {
    ok(existsSync("package-lock.json"),
        "package-lock.json should be present");
});

test("package-lock.json is included in published files", () => {
    const pkg = JSON.parse(readFileSync("package.json", "utf-8"));
    ok(pkg.files.includes("package-lock.json"),
        "package-lock.json should be in files array");
});
\end{lstlisting}
